\begin{minipage}{\TableWidth}
\small
\begin{tabularx}{\linewidth}{@{}L{1.25in}Y@{}}
\toprule
Part & Content \\
\midrule
Problem & The graphs of \(y=x^3-3x+2\) and \(x+4y=4\) intersect in the points \((x_1,y_1)\), \((x_2,y_2)\), and \((x_3,y_3)\). If \(x_1+x_2+x_3=A\) and \(y_1+y_2+y_3=B\), compute the ordered pair \((A,B)\). \\
\addlinespace
Noisy Prefix & From \(x+4y=4\), the prefix derives
\[
y=1-\frac{x}{4}.
\]
Substituting this into \(y=x^3-3x+2\) gives
\[
1-\frac{x}{4}=x^3-3x+2,
\]
and hence \(4x^3-11x+4=0\). The prefix then \textcolor{DenoiseRed}{attempts to find the individual roots using the Rational Root Theorem, testing candidates such as \(x=1\), \(x=\tfrac12\), \(x=-\tfrac12\), \(x=2\), and \(x=-1\)}. [omitted: several direct substitutions showing that these candidates are not roots.]

\textcolor{DenoiseRed}{The prefix keeps searching for explicit roots, although the problem only asks for the sums of the coordinates.} \\
\addlinespace
Model Continuation & \textbf{It seems the Rational Root Theorem isn't yielding roots easily. Let's try another method.} Given \(4x^3-11x+4=0\), the continuation writes it as
\[
4x^3+0x^2-11x+4=0.
\]
\textbf{By Vieta's formulas},
\[
x_1+x_2+x_3=-\frac{0}{4}=0,
\]
so \(A=0\). Using \(y=1-x/4\),
\[
y_1+y_2+y_3
=\left(1-\frac{x_1}{4}\right)+\left(1-\frac{x_2}{4}\right)+\left(1-\frac{x_3}{4}\right)
=3-\frac{x_1+x_2+x_3}{4}=3.
\]
Thus, the final answer is \(\boxed{(0,3)}\). \\
\bottomrule
\end{tabularx}
\end{minipage}
